% ---
% id: reductions/rem:precomposition-scope
% kind: rem
% title: "Scope of the precomposition lemma"
% summary: >-
%   The precomposition lemma needs only membership in the set the attack-set
%   displays cut out, at a separable input, and gives the inclusion
%   of the specified-source attack set, composed with any channel, in the
%   virtual-source one. A general entry has no reading beside a non-empty
%   source segment.
% prominence: standard
% status: accepted
% uses: [lem:precomposition]
% mentions: [def:attset, def:chain, thm:licence, thm:two-sided, rem:ordering, meat-core/def:admissible-declaration, transmitter/def:srcert]
% assumes: []
% refs: []
% provenance: newly-derived
% lean:
% ---

\begin{remark}\label{rem:precomposition-scope}
\emph{The general entry.} The right-hand side $\Theta$ of
Lemma~\ref{lem:precomposition}\,(iv), Hermitian and not a real multiple of the
identity, is the channel-level shape that the Hermitian right-hand side of a
\emph{general} entry asks for, an entry Definition~\ref{def:attset} records as
admissible input to the consuming contract and never licensed there. At the state level
such an entry beside a \tref{source link} fails earlier still, and this is
separate from the counterexample. There its right-hand side is a Hermitian
operator on Alice's retained system (Definition~\ref{def:attset};
\xref{meat-core/def:admissible-declaration}\,(vi)), and the retained system a
\tref{source-replacement certificate} is written against is the one of the
\emph{virtual} description, at every length of the \tref{source segment}
(\xref{transmitter/def:srcert}, ``One reading at every $L_{\mathrm s}$''); no
contract relates that system or its marginal to the specified description's.
So such an entry does not merely fail to be preserved: beside a non-empty
source segment it has no reading to preserve. That, and not the
counterexample alone, is why the admission rule of
Definition~\ref{def:chain} is a refusal.

\emph{Membership, and where it comes from.} The hypothesis of
Lemma~\ref{lem:precomposition}\,(i) to (iii) is \emph{membership} in the set
the displays of Definition~\ref{def:attset} cut out. The attack sets that
definition names are those of the \tref{round model}, whose input registers are
finite-dimensional; the displays themselves carry no condition on the
dimension of $E'$, and the lemma reads them at a separable input register,
so (i) to (iii) hold there.
Obtaining that membership from Theorem~\ref{thm:licence}\,(i), (ii) or (v) is
a different matter, because that theorem types its attack on a
finite-dimensional input register. Theorem~\ref{thm:two-sided}\,(ii)
therefore does not route the first round through the lemma: it applies
Theorem~\ref{thm:licence} to the pre-composed attack, whose input register is
the \tref{virtual transmitted system}. The lemma is the inclusion
Remark~\ref{rem:ordering} needs, and the second of the two routes in the
proof of Theorem~\ref{thm:two-sided}\,(ii).

\emph{The inclusion.} Write
$\AttSet^{\mathrm{real}}_1\coloneqq\AttSet_{\CList}(E',\cdot)$ for the set the
first-round attack of the specified-source protocol is admissible in and
$\AttSet^{\mathrm{virtual}}_1\coloneqq
\AttSet_{\CList}(E_{\mathrm{pre}},\cdot)$ for the set its precomposition with
$\Gamma$ must land in. Then Lemma~\ref{lem:precomposition}\,(i) to (iii) say
exactly
$\AttSet^{\mathrm{real}}_1\circ\Gamma\subseteq\AttSet^{\mathrm{virtual}}_1$,
\emph{for every channel $\Gamma$}. That is the inclusion
Remark~\ref{rem:ordering} names, and Theorem~\ref{thm:two-sided} is where it
is spent.
\end{remark}
